\section{Related Work}
\label{related_work}

\paragraph{Reachability via Linear Relaxation}
\looseness-1
While exact reachable set computation for nonlinear systems is typically intractable, reachable set overapproximations (RSOAs) can often be efficiently computed to certify safety \citep{DBLP:journals/arcras/AlthoffFG21}. When the closed-loop system involves NN components (e.g., learned dynamics or controller), RSOAs can be computed via linear relaxation-based NN verification (NNV) tools (e.g., CROWN \citep{Zhang2018CROWN}) \citep{Zhang2018CROWN, Everett2021NFL, Jafarpour2024AkashNFL}. Here, conservativeness grows with the number of nonlinearities \citep{sage-mp} and can be excessive for nonlinear robot dynamics. Symbolic and one-shot methods mitigate this but raise runtime and memory costs \citep{Chen2023One-Shot, Rober2024CARV, Everett2021NFL}. We instead model the robot with Koopman-linearized dynamics, reducing the number of nonlinearities in the CG and improving scalability (i.e., computation time, conservativeness).

\paragraph{Koopman Reachability}
\looseness-1
Koopman \citep{koopman_theory} and Carleman \citep{doi:10.1137/1.9781611976847.1, Bayer2023} linearizations accelerate reachability but introduce approximation error. Moreover, while error bounds exist, these derivations are limited to quadratic or polynomial dynamics \citep{DBLP:journals/pnas/LiuKKLTC21, DBLP:journals/corr/abs-1711-02552, doi:10.1137/1.9781611976847.1}, which do not apply to many robot dynamics models. Prior Koopman reachability methods also lack reachable set coverage calibration \citep{DBLP:journals/access/ThapliyalH22} or guarantee containment only for a finite set of observed trajectories \citep{DBLP:conf/cdc/KochdumperB22}, failing to ensure safety for previously-unseen trajectories of the \textit{true} system. Prior methods also do not support NN liftings and consider autonomous systems, which corresponds to fixing a single controller. Changing controllers, as often done in robotics, requires repeating the calibration \citep{DBLP:conf/rp/ForetsS21, DBLP:conf/adhs/BakBDGP21, DBLP:journals/fmsd/BakBHKLP25}, which is inefficient. In contrast, we support NN liftings, provide probabilistic coverage guarantees for the \textit{true} dynamics, and integrate control design, enabling data-efficient reachability across multiple controllers.

\paragraph{Data-Driven Reachability}
\looseness-1
Data-driven reachability uses black-box trajectory data. Some methods adapt Hamilton-Jacobi analysis \citep{reachblackboxdynamical}, but this is costly for high-dimensional systems. Others learn reachability functions \citep{neureach} but require dense data to accurately generalize, or yield ellipsoidal RSOAs \citep{10772635datadrivenreach} assuming privileged system data (e.g., Lipschitz constants). Sampling-based methods \citep{pmlr-v168-lew22a} are limited to short horizons. Scenario optimization provides probabilistic guarantees \citep{pmlr-v242-dietrich24a} but is slower than Koopman-based methods \citep{DBLP:journals/fmsd/BakBHKLP25}.
Moreover, these methods assume a fixed controller, yet robots often switch controllers to achieve different goals, altering the closed-loop dynamics and forcing reachable sets to be recomputed. Trajectory-tracking reachability provides reusable tracking-error sets across reference trajectories \citep{DBLP:conf/cdc/HerbertCHBFT17, DBLP:conf/wafr/SinghCHTP18, DBLP:journals/ral/KnuthCOB21}. Data-driven variants learn contraction metrics \citep{DBLP:conf/cdc/ChouOB21, DBLP:conf/wafr/ChouOB22, 10161001} for tracking-based reachability of unknown systems, but finding such metrics is generally difficult. We build on this work, yielding Koopman reachable sets that are fast to compute, scalable to long horizons, and reusable for various closed-loop dynamics induced by tracking different trajectories.